\documentclass[10pt,a4paper]{amsart}
\usepackage{amsmath,amssymb,amsthm,graphicx}
\usepackage{verbatim}
\newenvironment{reference}{\comment}{\endcomment}
%\newenvironment{reference}{}{}
\newenvironment{slogan}{\comment}{\endcomment}
\newenvironment{history}{\comment}{\endcomment}

% For commutative diagrams we use Xy-pic
\usepackage[all]{xy}

% We use 2cell for 2-commutative diagrams.
\xyoption{2cell}
% Theorem environments.
%
\theoremstyle{plain}
\newtheorem{theorem}[subsection]{Theorem}
\newtheorem{proposition}[subsection]{Proposition}
\newtheorem{lemma}[subsection]{Lemma}

\theoremstyle{definition}
\newtheorem{definition}[subsection]{Definition}
\newtheorem{example}[subsection]{Example}
\newtheorem{exercise}[subsection]{Exercise}
\newtheorem{situation}[subsection]{Situation}

\theoremstyle{remark}
\newtheorem{remark}[subsection]{Remark}
\newtheorem{remarks}[subsection]{Remarks}

\numberwithin{equation}{subsection}

% Macros
%
\def\lim{\mathop{\mathrm{lim}}\nolimits}
\def\colim{\mathop{\mathrm{colim}}\nolimits}
\def\Spec{\mathop{\mathrm{Spec}}}
\def\Hom{\mathop{\mathrm{Hom}}\nolimits}
\def\Ext{\mathop{\mathrm{Ext}}\nolimits}
\def\SheafHom{\mathop{\mathcal{H}\!\mathit{om}}\nolimits}
\def\SheafExt{\mathop{\mathcal{E}\!\mathit{xt}}\nolimits}
\def\Sch{\mathit{Sch}}
\def\Mor{\mathop{\mathrm{Mor}}\nolimits}
\def\Ob{\mathop{\mathrm{Ob}}\nolimits}
\def\Sh{\mathop{\mathit{Sh}}\nolimits}
\def\NL{\mathop{N\!L}\nolimits}
\def\CH{\mathop{\mathrm{CH}}\nolimits}
\def\proetale{{pro\text{-}\acute{e}tale}}
\def\etale{{\acute{e}tale}}
\def\QCoh{\mathit{QCoh}}
\def\Ker{\mathop{\mathrm{Ker}}}
\def\Im{\mathop{\mathrm{Im}}}
\def\Coker{\mathop{\mathrm{Coker}}}
\def\Coim{\mathop{\mathrm{Coim}}}

% Boxtimes
%
\DeclareMathSymbol{\boxtimes}{\mathbin}{AMSa}{"02}

%
% Macros for moduli stacks/spaces
%
\def\QCohstack{\mathcal{QC}\!\mathit{oh}}
\def\Cohstack{\mathcal{C}\!\mathit{oh}}
\def\Spacesstack{\mathcal{S}\!\mathit{paces}}
\def\Quotfunctor{\mathrm{Quot}}
\def\Hilbfunctor{\mathrm{Hilb}}
\def\Curvesstack{\mathcal{C}\!\mathit{urves}}
\def\Polarizedstack{\mathcal{P}\!\mathit{olarized}}
\def\Complexesstack{\mathcal{C}\!\mathit{omplexes}}
% \Pic is the operator that assigns to X its picard group, usage \Pic(X)
% \Picardstack_{X/B} denotes the Picard stack of X over B
% \Picardfunctor_{X/B} denotes the Picard functor of X over B
\def\Pic{\mathop{\mathrm{Pic}}\nolimits}
\def\Picardstack{\mathcal{P}\!\mathit{ic}}
\def\Picardfunctor{\mathrm{Pic}}
\def\Deformationcategory{\mathcal{D}\!\mathit{ef}}
\usepackage{cleveref}
\crefname{originallemma}{original lemma}{original lemmas}
\Crefname{originallemma}{Original lemma}{Original lemmas}
\makeatletter
\newlabel{spaces-properties-lemma-open-subspaces}{{03BZ}{source}{}{sourcepointer-1}{}}
\newlabel{spaces-properties-lemma-open-subspaces@cref}{{[originallemma][0][]03BZ}{[1][0][]source}{}{}{}}
\makeatother
\title{Two irreducible algebraic spaces are birational exactly when they share an isomorphic nonempty open subspace}
\author{The Stacks Project Authors}
\date{stacks-project:spaces-morphisms.tex:local-snapshot-20261012; GNU FDL 1.2 or later}
\begin{document}
\maketitle
\setcounter{subsection}{1}
\noindent\textit{Editorial note.} All delivered mathematical text is the author's own, copied from the pinned Stacks Project chapter file spaces-morphisms.tex; no statement or proof step was rewritten, reordered or removed.\par\medskip
\noindent\textit{Editorial note.} The proof closes with the author's sentence that the reader finds U' -> V' and V' -> U' to be mutually inverse; that sentence is kept verbatim. The construction of U' and V' and both factorizations are carried out inside the counted span, so only the routine verification of the two composites is left to the reader.\par\medskip
\noindent\textit{Editorial note.} The single reference of the counted spans is recorded as a non-load-bearing navigation pointer; the argument that consumes it stays inside the counted proof span.\par\medskip
\noindent\textit{Editorial note.} Preamble: the three native excerpt spans of sources/stacks-spaces-morphisms/preamble.tex (lines 8-13, 14-18 and 154-217). The chapter file itself carries no preamble text.\par\medskip
\noindent\textit{Allowed standard prerequisites (editorial).}\begin{itemize}
\item Dominant rational maps between algebraic spaces and the notion of birationality. used at proof lines 9897-9900. Supplies the two inverse dominant rational maps the argument starts from.
\item Every open subset of |U| that comes from an open of |V| is the underlying set of an open subspace. Properties of Spaces, Lemma spaces-properties-lemma-open-subspaces, used at proof lines 9908-9910. Produces the open subspace U' of U on which the two maps become inverse.
\item Irreducibility of |X| and |Y| makes dense opens nonempty and shrinkable. used at proof lines 9898-9907. Justifies the successive shrinkings of U and V.
\end{itemize}\medskip
\ifdefined\hypertarget\hypertarget{sourcepointer-1}{}\fi
\noindent\textit{Source reference.} Original lemma 03BZ: focused-seeds/AGP/sources/stacks-snapshot-20261011/tags.txt, source lines 4366--4366. Navigation pointer into the Stacks Project; the standard fact it supplies is applied inside the counted proof span and the argument resting on it stays there.\par\medskip
\begin{lemma}
\label{lemma-birational}
Let $S$ be a scheme. Let $X$ and $Y$ be algebraic
space over $S$ with $|X|$ and $|Y|$ irreducible.
Then $X$ and $Y$ are birational if and only if
there are nonempty open subspaces $U \subset X$ and $V \subset Y$
which are isomorphic as algebraic spaces over $S$.
\end{lemma}

\begin{proof}
Assume $X$ and $Y$ are birational. Let $f : U \to Y$ and $g : V \to X$
define inverse dominant rational maps from $X$ to $Y$ and from $Y$ to $X$.
After shrinking $U$ we may assume $f : U \to Y$ factors through $V$.
As $g \circ f$ is the identity as a dominant rational map, we see that
the composition $U \to V \to X$ is the identity on a dense open of $U$.
Thus after replacing $U$ by a smaller open we may assume that
$U \to V \to X$ is the inclusion of $U$ into $X$.
By symmetry we find there exists an open subspace $V' \subset V$
such that $g|_{V'} : V' \to X$ factors through $U \subset X$
and such that $V' \to U \to Y$ is the identity.
The inverse image of $|V'|$ by $|U| \to |V|$ is an open of $|U|$
and hence equal to $|U'|$ for some open subspace $U' \subset U$, see
Properties of Spaces, Lemma \ref{spaces-properties-lemma-open-subspaces}.
Then $U' \subset U \to V$ factors as $U' \to V'$.
Similarly $V' \to U$ factors as $V' \to U'$.
The reader finds that $U' \to V'$
and $V' \to U'$ are mutually inverse
morphisms of algebraic spaces over $S$
and the proof is complete.
\end{proof}

\end{document}
